\section{The elementary Gram background requires a structural correction (NS-70)}
\label{sec:ns70-background}

\paragraph{Scope.}
The endpoint model in NS-69 comes from the elementary part of Ehm's
complete $q=2$ Gram formula. Its scalar correction is small in some
entries, but that does not make it a small quadratic-form perturbation.
We first identify and repair a mismatch in diagonal regularity. Even
after that repair, a uniform comparison on all coefficient vectors is
excluded. None of these statements is a lower bound or obstruction
for the actual optimized residual direction.

For positive real $u,v$ let $G(u,v)=\langle\sigma_u,\sigma_v\rangle$.
With $K_1,K_2$ from Section~\ref{sec:ns61-gram-remainder}, define
\[
 K(u,v)=\frac{K_2+K_1|\log(u/v)|+\tfrac14\log^2(u/v)}{\max(u,v)},
 \qquad S(u,v)=G(u,v)-K(u,v).
\]
At integer indices this is the full arithmetic remainder of the Ehm
formula. Set $d=K_2-2K_1$. The constants are certified below; in
particular $d>0$ and $K_1>3/2$.

\begin{proposition}[A small entrywise correction cancels the diagonal cusp]
\label{prop:ns70-cusp}
Both $G$ and $K$ are positive definite on every finite collection of
distinct positive dilations. For the adjacent coefficient vector
corresponding to
$\epsilon_n=\sigma_n-(n-1)\sigma_{n-1}/n$, write $G[\epsilon_n]$
and $K[\epsilon_n]$ for its two quadratic forms. Then
\begin{equation}
 G[\epsilon_n]\sim\frac\kappa{n^3},\qquad
 K[\epsilon_n]\sim\frac d{n^2},\qquad
 \frac{S[\epsilon_n]}{K[\epsilon_n]}\longrightarrow-1.
 \label{eq:ns70-cusp-asymptotic}
\end{equation}
Consequently no fixed $c>0$ can satisfy $G\geq cK$ on all finite
integer coefficient vectors. In particular, no relative remainder
bound $|S[a]|\leq\delta K[a]$ with a fixed $\delta<1$ can hold
on that class. This follows from local regularity and does not use
an assumption about the location of zeta zeros.
\end{proposition}
\begin{proof}
Write $r=\log(v/u)\geq0$. Normalizing by $\sqrt{uv}$ turns $K$
into the translation kernel
\[
 \phi(r)=e^{-|r|/2}(K_2+K_1|r|+r^2/4).
\]
Its Fourier density, with
$\phi(r)=(2\pi)^{-1}\int\widehat\phi(t)e^{itr}dt$, is
\begin{equation}
 w_K(t)=\frac{d t^4+(K_2/2-3/2)t^2
                         +K_2/16+K_1/8+1/8}{(t^2+1/4)^3}>0.
 \label{eq:ns70-background-symbol}
\end{equation}
This is obtained by differentiating
$\int e^{-a|r|}e^{-itr}dr=2a/(a^2+t^2)$ twice in $a$ and then
putting $a=1/2$. The three numerator coefficients are positive by
the recorded constant enclosures. Hence $K$ is strictly positive
definite; $G$ is already a Gram of independent dilations.

For $h_n=\log(n/(n-1))$ direct evaluation gives
\begin{equation}
 n^2K[\epsilon_n]
 =K_2-2K_1(n-1)h_n-\tfrac12(n-1)h_n^2\longrightarrow d.
 \label{eq:ns70-exact-adjacent-background}
\end{equation}
NS-61 gives $n^3G[\epsilon_n]\to\kappa$ and its complete upper
bound. Thus the ratio $G[\epsilon_n]/K[\epsilon_n]$ tends to zero,
proving all assertions. The cusp is explicit:
$\phi'(0^+)=K_1-K_2/2=-d/2$, while the actual smoothed Gram
has the additional order of cancellation in
\eqref{eq:ns70-cusp-asymptotic}.
\end{proof}

\paragraph{The cusp correction is exactly diagonal in adjacent differences.}
For $B(u,v)=1/\max(u,v)$, direct evaluation gives
\[
 B[\epsilon_m,\epsilon_n]=\frac{\mathbf1_{m=n}}{n^2},
 \qquad B[\sigma_j,\epsilon_n]=0\quad(j<n).
\]
Indeed for $j<n$ the right difference is
$1/n-((n-1)/n)/(n-1)=0$, including $j=n-1$.
The diagonal difference equals $1/n^2$. Consequently the repair
subtracts exactly $d\,\operatorname{diag}(n^{-2})$ from the raw
new-difference Gram and leaves every elementary old/new cross column
unchanged. In particular it cannot improve the endpoint model $L$
in NS-69: the correction cancels from those correlations identically.
Its role is to repair the energy assigned to correction directions.
The old model Gram also changes, so a projected model must still
recompute its complete Schur complement.

\begin{proposition}[Repairing the cusp still leaves the zeta-weight comparison]
\label{prop:ns70-repaired-background}
The corrected elementary kernel
\begin{equation}
 K_0(u,v)=K(u,v)-\frac d{\max(u,v)}
 =\frac{2K_1+K_1|\log(u/v)|+\tfrac14\log^2(u/v)}{\max(u,v)}
 \label{eq:ns70-repair}
\end{equation}
is strictly positive definite and has no first-derivative cusp.
It satisfies
\[
 K_0[\epsilon_n]\sim\frac{K_1-1/2}{n^3},\qquad
 \frac{G[\epsilon_n]}{K_0[\epsilon_n]}
                \longrightarrow\frac\kappa{K_1-1/2}.
\]
Nevertheless there are no constants $c>0$ or $C<\infty$ such that,
on all finite integer coefficient vectors, respectively
\[
 G\geq cK_0\qquad\text{or}\qquad G\leq CK_0.
\]
Both failures are unconditional. They exclude uniform all-coefficient
comparisons, not selected-direction estimates, finite-block inequalities
with size-dependent constants, or the optimized NB criterion.
\end{proposition}
\begin{proof}
The Fourier density of the repaired kernel is
\begin{equation}
 w_0(t)=\frac{(K_1-3/2)t^2+K_1/4+1/8}{(t^2+1/4)^3}>0.
 \label{eq:ns70-repaired-symbol}
\end{equation}
This proves strict positivity and the cancellation of the cusp.
In~\eqref{eq:ns70-exact-adjacent-background} replace $K_2$ by
$2K_1$ and use
$(n-1)h_n=1-1/(2(n-1))+O(n^{-2})$ and
$(n-1)h_n^2=1/(n-1)+O(n^{-2})$.
This gives the displayed adjacent asymptotic.

The complete Gram has Mellin density
\[
 w_G(t)=\frac{|\zeta(1/2+it)|^2}{(t^2+1/4)^2}.
\]
For clarity, a uniform quadratic-form inequality on all finite integer
coefficients would force the corresponding pointwise inequality between
$w_G$ and $w_0$. Here is the full localization argument. Both densities
are continuous, even and integrable. Their homogeneous kernels extend
to all positive real dilation parameters. Replace any finite set of real
parameters $u_j$ by integers nearest to $D u_j$ and multiply both forms
by $D$; kernel continuity lets $D\to\infty$. Thus an integer inequality
extends to all real parameters. Riemann sums then extend it to compactly
supported real coefficient densities in $r=\log u$. Choose long real
cosine packets in $r$. Their squared Fourier transforms form approximate
identities at the paired frequencies $\pm t_0$. Continuity near the
frequencies and integrability away from them give the pointwise inequality.
This reasoning retains all quadratic cross terms and uses arbitrary
finite real coefficients, not an arithmetic coefficient rule.

A known critical-line zero makes $w_G(t_0)=0$ while $w_0(t_0)>0$,
so the lower comparison is impossible. For the upper comparison,
\[
 \frac{w_G(t)}{w_0(t)}
 =|\zeta(1/2+it)|^2\,
   \frac{t^2+1/4}{(K_1-3/2)t^2+K_1/4+1/8}.
\]
The rational factor tends to the positive constant $(K_1-3/2)^{-1}$.
Known unbounded critical-line values, for example
\cite[Theorem 1]{Soundararajan2008}, therefore exclude any finite $C$.
The lower failure uses zeros on the line, not hypothetical off-line zeros.
Neither failure says that the actual Gram, or an original Weil form, has
negative energy.
\end{proof}

\paragraph{Complete constant and adjacent-vector certificates.}
Arb at 256 and 384 bits encloses every numerator coefficient in
\eqref{eq:ns70-background-symbol} and
\eqref{eq:ns70-repaired-symbol} strictly above zero, with
$.0091604<d<.0091605$. Complete two-vector Gram comparisons give
\begin{center}\small
\begin{tabular}{rcc}
\toprule
$n$ & $G[\epsilon_n]/K[\epsilon_n]$
    & $G[\epsilon_n]/K_0[\epsilon_n]$\\
\midrule
256 & $(.36207,.36209)$ & $(1.11185,1.11187)$\\
1024 & $(.11987,.11989)$ & $(1.11423,1.11425)$\\
4096 & $(.03260,.03262)$ & $(1.11498,1.11500)$\\
\bottomrule
\end{tabular}\end{center}
The true quadratic energy stays positive. The Bernoulli/Hurwitz
remainder is retained; a doubled series cutoff gives overlapping
values for all $34$ constants and quantities. The asymptotic assertions
follow from the proofs, not from this finite table.

\paragraph{What had to change.}
The elementary background had the wrong local regularity; removing
$d/\max(u,v)$ repairs that mismatch and recovers the correct adjacent
power $n^{-3}$. This repair is real but insufficient: the complete
arithmetic factor is still neither uniformly above nor uniformly below
the repaired background. A useful next bound must keep the particular
residual direction or supply an additional arithmetic estimate. The
finite endpoint remainder in NS-69 is not silently discarded.
